% Section 1.7: PAC Generation from SVD Files

\section{Generating Peripheral Access Crates (PAC)}
\label{sec:pac-generation}

The journey from vendor hardware specifications to usable Rust code begins with System View Description (SVD) files. These XML files describe every register, field, and peripheral in the microcontroller. In this section, we'll explore how to generate a Peripheral Access Crate (PAC) from these vendor-provided files.

\subsection{What is an SVD File?}
\label{subsec:svd-file}

System View Description (SVD) files are XML documents provided by microcontroller vendors that describe the complete memory-mapped register layout of a chip. They contain:

\begin{itemize}
    \item Register addresses and sizes
    \item Field definitions (bit positions, access types)
    \item Peripheral descriptions
    \item Interrupt definitions
    \item Memory region information
\end{itemize}

\begin{notebox}
    SVD files are the "single source of truth" for register definitions. Any discrepancy between the SVD and the actual hardware requires manual verification against the reference manual.
\end{notebox}

\subsection{Obtaining SVD Files}
\label{subsec:obtaining-svd}

For the EFR32MG24, SVD files are available through:

\begin{enumeratex}
    \item Download the Silicon Labs Gecko Platform DFP (Device Family Pack)
    \item Extract the SVD file for your specific variant
    \item Verify the variant matches your hardware (e.g., EFR32MG24B220F1536IM48)
\end{enumeratex}

\begin{minted}{bash}
# Download DFP pack
wget https://www.silabs.com/documents/public/cmsis-packs/\
    SiliconLabs.GeckoPlatform_EFR32MG24_DFP.2025.6.2.pack

# Extract SVD
unzip SiliconLabs.GeckoPlatform_EFR32MG24_DFP.2025.6.2.pack \
    "SVD/EFR32MG24/*.svd" -d svd/
\end{minted}

\subsection{Critical Discovery: A020 vs B220}
\label{subsec:svd-variant-discovery}

Our initial testing revealed a critical lesson: not all EFR32MG24 variants are identical. We initially used the A020 SVD (first available) but our hardware is the B220 variant with different register addresses.

\begin{warning}
    Generating a PAC from the wrong SVD variant produces code that compiles but fails on hardware, leading to subtle bugs that are difficult to debug.
\end{warning}

\subsection{Using svd2rust}
\label{subsec:using-svd2rust}

\tool{svd2rust} converts SVD files into Rust code:

\begin{minted}{bash}
cd efr32mg24-pac
svd2rust -i ../svd/EFR32MG24B220F1536IM48.svd \
         --target cortex-m \
         -o src/
\end{minted}

\subsection{Understanding PAC Output}
\label{subsec:pac-output}

The generated PAC contains:

\begin{itemize}
    \item 133 peripherals
    \item 138,448 lines of Rust code
    \item Register-level access with type safety
    \item Critical section protection
\end{itemize}

\subsection{Register Access Patterns}
\label{subsec:register-access}

With the PAC, you can access hardware registers directly:

\begin{exampleoutput}
PAC access looks like this:
\end{exampleoutput}

\begin{minted}{rust}
use efr32mg24_pac::GPIO_S as GPIO;

// Read-modify-write atomic operation
critical_section::with(|_| {
    GPIO.porta_doutset.write(|w|
        w.doutset().bits(1 << 5)  // Set PA5 high
    );
});
\end{minted}

\begin{keytakeaways}
    \item Always verify SVD variant matches your hardware
    \item Never edit PAC files manually—regenerate if needed
    \item Use critical sections for atomic register operations
    \item PAC provides zero-cost abstractions over raw memory
\end{keytakeaways}

\begin{exercises}
    \item\textbf{Easy}: Download the DFP pack for your specific variant and extract the SVD file.
    
    \item\textbf{Medium}: Generate a PAC using svd2rust and verify it compiles successfully.
    
    \item\textbf{Hard}: Compare register addresses between A020 and B220 SVD files and identify at least three differences.
\end{exercises}
